% Sample book for the style tests. It exercises every part of the style once:
% front matter, a part page, a chapter opener with a long title, running heads
% with a long section title, a figure, a table, a code example, the callouts,
% an index entry and the colophon.
\documentclass[10pt,twoside,openright]{book}
\usepackage{style/animalbook}
\input{style/options.tex}

\newcommand{\filler}{A probe asks whether a property can be read from the model's
internal state; steering asks whether pushing the state along a direction changes
what the model does. The first is about information, the second about cause, and
most confusions come from mixing the two. A difference of means is the simplest
probe: average the positives, average the negatives and take the difference.\par}

\begin{document}
\frontmatter
\halftitlepage
\copyrightpage{3 November 2026}{Built by the test suite of the book skill. Private:
not for circulation.}
\tableofcontents

\chapter{Conventions Used in This Book}

\workedexample{title}{A number from the project, computed step by step.}

\mainmatter
\part{Foundations}

\chapter{Probes and Steering}
\label{ch:probes}

\sources{The probe code in \code{src/sample\_project/dom.py} and
\code{scripts/landmark\_probe.py}; results in \code{outputs/probes/} as cited in the
handouts. Meetings only for why we did each step.}

\filler

\keyidea{An SAE rewrites one activation vector as a short list of (feature, value)
pairs.}

\careful{A probe that works shows that the property is linearly readable from the
state. It does not show that the model uses that direction.}

\workedexample{nothing stands out on 10 episodes}{The largest null value is 5.82, so
no feature passes the threshold \ev{negative}.}

\inproject{We save the exact token IDs of every model call, then replay them through
the same pinned model with a hook that saves the residual stream.}

\questionsfor{the supervisor}{Is the detection question parked for good, or should the
midterm still report a detection baseline?}

\todo{P2}{Explain the logistic probe and its regularization.}

\smallnote{Source: \code{outputs/run\_01/probe\_metrics.json}.}

\filler

\section{What a Probe Is}

Take one activation vector $x \in \mathbb{R}^{d}$ per example. A \textbf{linear
probe}\index{probe} is a weight vector $w$ and a bias $b$; its score is
\[
  s(x) = w^\top x + b .
\]
\filler\filler\filler

\section{A Section Title That Is Long Enough to Run Past the Width of the Running Head}

\filler\filler

\begin{figure}[htbp]
  \centering
  \rule{0.6\linewidth}{0.25\linewidth}
  \caption{A placeholder figure with a caption long enough to wrap onto a second line
  of text under it}
  \label{fig:placeholder}
\end{figure}

\begin{table}[htbp]
  \caption{Probe families used in this book}
  \centering
  \begin{tabular}{lll}
    \toprule
    Probe & Training & Used in \\
    \midrule
    DoM & none & Chapter~\ref{ch:probes} \\
    Logistic & L2, 5-fold & Chapter~\ref{ch:long} \\
    \bottomrule
  \end{tabular}
\end{table}

\begin{codeexample}{A DoM direction and its AUROC}
from sample_project.dom import auroc

direction = pos.mean(0) - neg.mean(0)
score = auroc(pos @ direction, neg @ direction)
\end{codeexample}

\filler\filler\filler\filler

\chapter{A Chapter Title That Is Long Enough to Need Two Lines on Its Opener}
\label{ch:long}

\filler\filler\filler

\subsection{A Subsection}

\filler\filler\filler\filler\filler

\backmatter
\printindex

\begin{colophon}
The animal on the cover of \emph{The Sample Book} is the common cuckoo
(\emph{Cuculus canorus}), a bird that lays its eggs in other birds' nests and
leaves the hosts to raise its young.
\end{colophon}

\end{document}
